\documentclass{article}

\usepackage{amsmath,amssymb}
\usepackage{xurl}
\usepackage[colorlinks=true,linkcolor=blue,citecolor=blue,urlcolor=blue]{hyperref}

% Shared commands for notes in docs/paper-gaps/.

\DeclareUnicodeCharacter{2080}{\ifmmode{}_{0}\else\textsubscript{0}\fi}
\DeclareUnicodeCharacter{2081}{\ifmmode{}_{1}\else\textsubscript{1}\fi}
\DeclareUnicodeCharacter{2082}{\ifmmode{}_{2}\else\textsubscript{2}\fi}
\DeclareUnicodeCharacter{2099}{\ifmmode{}_{n}\else\textsubscript{n}\fi}

\newcommand{\Lean}{\textsc{Lean}}
\ExplSyntaxOn
\NewDocumentCommand{\leanid}{m}{
  \ifmmode
    \textsf{\detokenize{#1}}
  \else
    \group_begin:
    \urlstyle{sf}
    \tl_set:Nn \l_tmpa_tl {#1}
    \tl_replace_all:Nnn \l_tmpa_tl {\_} {_}
    \exp_args:NV \path \l_tmpa_tl
    \group_end:
  \fi
}
\ExplSyntaxOff
\providecommand{\tr}{\operatorname{tr}}
\newcommand{\tnleanIssueBase}{https://github.com/LionSR/TNLean/issues/}
\newcommand{\tnleanPrBase}{https://github.com/LionSR/TNLean/pull/}
\newcommand{\tnleanDocBase}{https://sirui-lu.com/TNLean/docs/}
\newcommand{\ghissue}[1]{\href{\tnleanIssueBase#1}{GitHub issue \##1}}
\newcommand{\ghpr}[1]{\href{\tnleanPrBase#1}{PR \##1}}
\newcommand{\doclink}[1]{\href{\tnleanDocBase#1.html}{\path{#1}}}

% Dirac notation and the identity operator, as in blueprint/src/macros/common.tex.
% \providecommand keeps note-local definitions authoritative.
\usepackage{dsfont}
\providecommand{\ket}[1]{|#1\rangle}
\providecommand{\bra}[1]{\langle #1|}
\providecommand{\braket}[2]{\langle #1 | #2 \rangle}
\providecommand{\ketbra}[2]{|#1\rangle\!\langle #2|}
\providecommand{\Id}{\mathds{1}}
\providecommand{\one}{\mathds{1}}
\providecommand{\C}{\mathbb{C}}
\providecommand{\MN}[1]{M_{#1}(\C)}
\providecommand{\MD}{M_D(\C)}

% External referencing for paper sources.  The build helper writes sanitized
% auxiliary files under build/paper-aux/ and excludes duplicated source labels.
\usepackage{xr}

% Import a generated paper auxiliary file with a caller-supplied label prefix.
% The candidate paths support builds from docs/paper-gaps/, the repository root,
% and build/paper-aux/ itself.
\newcommand{\importpaperaux}[2]{%
  \IfFileExists{../../build/paper-aux/#2.aux}{%
    \externaldocument[#1]{../../build/paper-aux/#2}%
  }{%
    \IfFileExists{build/paper-aux/#2.aux}{%
      \externaldocument[#1]{build/paper-aux/#2}%
    }{%
      \IfFileExists{#2.aux}{%
        \externaldocument[#1]{#2}%
      }{}%
    }%
  }%
}

% General external reference macros
\newcommand{\paperref}[2]{\ref{#1-#2}}
\newcommand{\papereqref}[2]{(\ref{#1-#2})}

% CPSV16 source (1606.00608 / MPDO-22-12-17-2.tex) external document and shortcuts
\importpaperaux{cpsv16-}{1606.00608/MPDO-22-12-17-2-tnlean}
\newcommand{\cpsvref}[1]{\paperref{cpsv16}{#1}}
\newcommand{\cpsveqref}[1]{\papereqref{cpsv16}{#1}}

% Verdict marker: \gapnote{<kind>}{<status>}. Kind is the policy
% classification (clarification, local-correction, scope-restriction,
% unfaithful, false-source, open-gap); status is open, wip (elimination
% underway), resolved, or historical. Machine-read by the site index;
% prints a verdict line.
\newcommand{\gapnote}[2]{\par\noindent\textbf{Verdict:} #1 (#2).\par}


\title{A Two-Site Counterexample to Cyclic MPS Bond Uniformity}
\author{}
\date{2026-08-31}

\begin{document}
\maketitle
\gapnote{false-source}{resolved}

\section{The printed corollary}

The Applications section of Ref.~\cite{Molnar2018NormalPEPS} considers a
closed matrix product state whose local matrices have sizes
$D_k\times D_{k+1}$.  Its translational-invariance corollary asserts that
translational invariance of the state and injectivity of the local tensors
imply equality of all virtual dimensions and the existence of one repeated
tensor.  For this finite closed chain, the source writes translational
invariance as equality with the one-site cyclic shift.
Thus its bond-dimension conclusion is
\begin{align}
    D_1=D_2=\cdots=D_n.
    \label{eq:peps_cyclic_mps_two_site_bond_uniformity_printed_conclusion}
\end{align}
The corollary itself does not state a lower bound on $n$.\footnote{Local source:
\path{Papers/1804.04964/paper_normal.tex:1803--1805}; the cyclic-shift equality
is displayed at lines 1807--1827.}

The surrounding MPS section of Ref.~\cite{Molnar2018NormalPEPS} says that its
Fundamental Theorem concerns chains with at least three sites, and the theorem
invoked in the corollary is stated under that condition.\footnote{Local source:
\path{Papers/1804.04964/paper_normal.tex:145} and
\path{Papers/1804.04964/paper_normal.tex:688--725}.}  The omission in the
corollary is therefore consequential rather than a harmless simplification:
the literal two-site specialization is false.

\section{A positive-dimensional two-site example}

Let the physical space be $\mathbb C^2$, with standard basis $e_0,e_1$, and
take
\begin{align}
    D_1&=1,
    &D_2&=2,
    \label{eq:peps_cyclic_mps_two_site_bond_uniformity_dimensions}\\
    A_1^i&=e_i^{\mathsf T}\in\operatorname{Mat}_{1\times2}(\mathbb C),
    &A_2^i&=e_i\in\operatorname{Mat}_{2\times1}(\mathbb C)
    \qquad (i=0,1).
    \label{eq:peps_cyclic_mps_two_site_bond_uniformity_tensors}
\end{align}
The first local tensor sends the two virtual basis vectors of
$\mathbb C^1\otimes\mathbb C^2$ to $e_0,e_1$; the second does the same from
$\mathbb C^2\otimes\mathbb C^1$.  Both local tensor maps are therefore
isomorphisms, in particular injective.

For physical indices $i,j\in\{0,1\}$, the coefficient before and after the
cyclic shift is
\begin{align}
    \tr(A_1^iA_2^j)
        &=e_i^{\mathsf T}e_j
         =\delta_{ij},
    \label{eq:peps_cyclic_mps_two_site_bond_uniformity_original_coefficient}\\
    \tr(A_2^iA_1^j)
        &=\tr(e_i e_j^{\mathsf T})
         =\delta_{ij}.
    \label{eq:peps_cyclic_mps_two_site_bond_uniformity_shifted_coefficient}
\end{align}
Equations~(\ref{eq:peps_cyclic_mps_two_site_bond_uniformity_original_coefficient})
and~(\ref{eq:peps_cyclic_mps_two_site_bond_uniformity_shifted_coefficient})
show that the two-site state is cyclically invariant.  Every virtual space is
nonzero and both local tensors are injective, while
(\ref{eq:peps_cyclic_mps_two_site_bond_uniformity_dimensions}) contradicts
(\ref{eq:peps_cyclic_mps_two_site_bond_uniformity_printed_conclusion}).

The construction and the negation of the literal two-site statement have
kernel-checked proofs.\footnote{Formal declaration:
\leanid{TNLean.PEPS.CycleMPSTwoSiteCounterexample.twoSiteBondUniformityStatement_false}
in \doclink{TNLean/PEPS/CycleMPSTwoSiteCounterexample}.}

\section{Corrected scope}

The source proof compares the site-dependent chain with its cyclic shift by
the injective MPS Fundamental Theorem, writes every $A_i$ in terms of $A_1$
and the matrices $L_i,R_i$, and then telescopes these matrices to construct a
repeated tensor $B$.  That argument is available in the source's stated regime
$n\geq3$ and does not apply to the example above.  The corrected statement
therefore retains the standing three-site hypothesis:
\begin{align}
    n\geq3
    \quad\Longrightarrow\quad
    D_1=D_2=\cdots=D_n,
    \label{eq:peps_cyclic_mps_two_site_bond_uniformity_corrected_scope}
\end{align}
together with the positive-dimensional convention for the virtual spaces.
This is precisely the regime in which the graph-theoretic cycle has one edge
for each bond and the cited Fundamental Theorem comparison is valid.

The complete corrected conclusion, including equality of all bond dimensions
and the repeated injective tensor $B$, is formalized by
\leanid{TNLean.PEPS.exists_constant_injectiveMPS_of_cyclicShiftInvariantState_per_bond}.
It is recorded as complete in the Blueprint node
\path{cor:exists_constant_injectiveMPS_of_cyclicShiftInvariantState_per_bond_pos}.

\section{Verdict}

The translational-invariance corollary of Ref.~\cite{Molnar2018NormalPEPS} is
false as printed if it is read at two sites.  The example
(\ref{eq:peps_cyclic_mps_two_site_bond_uniformity_tensors}) satisfies the
nondegenerate source hypotheses and has a cyclically invariant state, but its
two virtual dimensions are unequal.  The corrected statement is the
at-least-three-site result
(\ref{eq:peps_cyclic_mps_two_site_bond_uniformity_corrected_scope}), in
accordance with the standing scope and the Fundamental Theorem used by the
source proof.  The impossible unqualified reading is therefore a resolved
false-source statement, not a separate open proof obligation.

\bibliographystyle{alpha}
\bibliography{references}

\end{document}
